\begin{helper}[The full two-color component has a certified last exit]
In the data $G,k,e_0,\varphi,T,C,D,R$, let $\alpha$ be present at every
vertex of $X=\noderef{ScheideEnlargedVertexSet}(R)$.  Then there exists a
\noderef{ScheideLastExitPath} certificate for $\alpha$.
\end{helper}
\begin{proof}
Write $y_0=C.y(0)$ and $\alpha_0=D.\alpha_0$.  The extension field of
\noderef{ScheideReducedConfiguration} supplies a
\noderef{ScheideStructuralConfiguration}; its base-position field and the image
definition \noderef{TreeSequenceVertexSet} put $y_0$ in $V(T)$, hence in
$X$.  The setup \noderef{ScheideDefectiveSetup} gives
$\alpha_0\in\noderef{MissingColors}(G,\varphi,y_0)$.  Since $\alpha$ is
present at $y_0$, the complementary definitions of
\noderef{PresentColors} and \noderef{MissingColors} give
$\alpha\ne\alpha_0$.

The globally maximum field of that
\noderef{ScheideStructuralConfiguration}, expanded through
\noderef{GlobalMaximumTriple} and \noderef{KTriple}, says that $\varphi$
is proper.  At $y_0$ there is no incident $\alpha_0$-edge and there is
at least one incident $\alpha$-edge.  Properness permits at most one such
$\alpha$-edge: two distinct identities incident with $y_0$ would be
\noderef{AdjacentEdges} of the same color.  Consequently the vertex-node
$y_0$ has degree one in the two-color incidence graph.

Apply \noderef{TwoColorComponentSpanningPath}.  It supplies a simple path
whose support is the full component through $y_0$.  By
\noderef{TwoColorIncidenceDegreeAtMostTwo}, no internal node of this path
can have another component incidence.  Since $y_0$ has degree one, it is
one of the two path ends; reverse the path if necessary so that it starts
at $y_0$.  Its other end is a vertex-node.  Indeed, an edge-node in the
component has the two distinct endpoint vertex-neighbors specified by
\noderef{Incident} and the loopless law in
\noderef{LooplessMultigraph}, and hence cannot be a degree-one end.
Denote the other endpoint vertex by $t$.

The set $X$ is elementary by \noderef{ScheideEnlargedElementary}.  If
$t\in X$, then $\alpha$ is present at $t$.  Since $t$ is a degree-one
end of the alternating path, its unique incident path edge has color
$\alpha$ and $\alpha_0$ is missing there.  Thus the distinct vertices
$y_0,t\in X$ both miss $\alpha_0$, contrary to
\noderef{ColoredVertexSetElementary}.  (They are distinct because the
component path has a degree-one end $y_0$ and contains its incident
$\alpha$-edge.)  Hence $t\notin X$.

The finite path begins at a vertex of $X$ and ends at a vertex outside
$X$.  Choose the greatest path index $i$ whose node is a vertex-node in
$X$, and call that vertex $v$.  The bipartite definition
\noderef{TwoColorIncidenceGraph} makes the next two nodes an edge-node
$g$ and a vertex-node $u$.  Maximality of $i$ gives $u\notin X$ and says
that every later vertex-node is outside $X$.  The two adjacencies give
that $g$ is incident with $v$ and $u$, so
$g\in\noderef{BoundaryEdges}(G,X)$.  Its color is $\alpha_0$ or
$\alpha$.  But $X$ is closed by \noderef{ScheideEnlargedClosed}; since
$y_0\in X$ misses $\alpha_0$, the definition
\noderef{ColoredVertexSetClosed} forbids an $\alpha_0$-colored boundary
edge.  Therefore $\varphi(g)=\alpha$.  Finally $g$ is a node of the path,
whose support is the component, so its edge-node belongs to that
component.  These choices and facts are exactly all fields of the
claimed \noderef{ScheideLastExitPath} certificate.
\end{proof}
